\clearpage\pagestyle{otMain}\setlength{\otBlocEspacement}{2.5mm}\fontsize{11}{15}\selectfont\corrigesChapitre{15}{Le cosinus dans le triangle rectangle}
\begin{blocLivre}{cor-15-1}\begin{corrige}{1}{Nommer les côtés}\demarche{Repérer d’abord l’angle droit, puis l’angle étudié.}
L’hypoténuse est $\res{BC}$, opposée à l’angle droit en A. Pour l’angle B, le côté adjacent qui n’est pas l’hypoténuse est $\res{AB}$. AC est opposé à l’angle B.
\label{fig-cor-15-1}\figureGeo[corrige]{cos-01}
\end{corrige}\end{blocLivre}
\begin{blocLivre}{cor-15-2}\begin{corrige}{2}{Calculer un cosinus}\demarche{Écrire adjacent sur hypoténuse.}
Le côté adjacent à B est AB et l’hypoténuse BC. Donc $\cos\widehat B=6/10=\res{3/5=0{,}6}$. Le quotient est bien entre 0 et 1.
\label{fig-cor-15-2}\figureGeo[corrige]{cos-02}
\end{corrige}\end{blocLivre}
\begin{blocLivre}{cor-15-3}\begin{corrige}{3}{Changer d’angle}\demarche{Réidentifier l’adjacent pour chaque angle.}
Pour C, le côté adjacent est AC. Ainsi $\cos\widehat C=AC/BC=8/10=\res{0{,}8}$. Le côté adjacent dépend de l’angle choisi.
\label{fig-cor-15-3}\figureGeo[corrige]{cos-03}
\end{corrige}\end{blocLivre}
\begin{blocLivre}{cor-15-4}\begin{corrige}{4}{Un angle de soixante degrés}\demarche{Isoler le numérateur du rapport.}
$AB/12=\cos60^\circ=0{,}5$. Donc $AB=12\times0{,}5=\res6$ cm. On multiplie l’hypoténuse par le cosinus pour obtenir l’adjacent.
\label{fig-cor-15-4}\figureGeo[corrige]{cos-04}
\end{corrige}\end{blocLivre}
\begin{blocLivre}{cor-15-5}\begin{corrige}{5}{Une longueur approchée}\demarche{Multiplier par le cosinus en mode degrés.}
$AB=8\cos35^\circ\approx\res{6{,}55\ \mathrm{cm}}$. Le côté adjacent est plus petit que l’hypoténuse 8 cm, ce qui contrôle l’ordre de grandeur.
\label{fig-cor-15-5}\figureGeo[corrige]{cos-05}
\end{corrige}\end{blocLivre}
\begin{blocLivre}{cor-15-6}\begin{corrige}{6}{Retrouver l’hypoténuse}\demarche{Diviser l’adjacent par le cosinus.}
$\cos40^\circ=7/BC$. Donc $BC=7/\cos40^\circ\approx\res{9{,}14\ \mathrm{cm}}$. BC doit dépasser 7 cm ; multiplier par le cosinus aurait donné une longueur trop petite.
\label{fig-cor-15-6}\figureGeo[corrige]{cos-06}
\end{corrige}\end{blocLivre}
\begin{blocLivre}{cor-15-7}\begin{corrige}{7}{Déterminer un angle}\demarche{Utiliser le cosinus inverse du rapport exact.}
$\cos\widehat B=5/13$. La commande cosinus inverse en degrés donne $\widehat B\approx67{,}38^\circ$, soit $\res{67^\circ}$ au degré près.
\label{fig-cor-15-7}\figureGeo[corrige]{cos-07}
\end{corrige}\end{blocLivre}
\begin{blocLivre}{cor-15-8}\begin{corrige}{8}{Un angle puis l’autre}\demarche{Utiliser le complément à 90 degrés.}
$\widehat B=\arccos(9/15)\approx\res{53{,}1^\circ}$. Les deux angles aigus sont complémentaires : $\widehat C=90^\circ-\widehat B\approx\res{36{,}9^\circ}$. Utiliser l’angle non arrondi pour le second calcul.
\label{fig-cor-15-8}\figureGeo[corrige]{cos-08}
\end{corrige}\end{blocLivre}
\begin{blocLivre}{cor-15-9}\begin{corrige}{9}{L’adjacent à C}\demarche{Associer le bon petit côté à l’angle donné.}
Pour l’angle C, AC est l’adjacent. $AC=10\cos50^\circ\approx\res{6{,}43\ \mathrm{cm}}$. L’autre petit côté ne figure pas dans ce rapport.
\label{fig-cor-15-9}\figureGeo[corrige]{cos-09}
\end{corrige}\end{blocLivre}
\begin{blocLivre}{cor-15-10}\begin{corrige}{10}{Une hypothèse manquante}\demarche{Contrôler les conditions du théorème.}
Non : la formule adjacent sur hypoténuse de ce chapitre suppose un triangle rectangle, ici en A. Deux longueurs seules ne donnent pas cette hypothèse. Il faut la connaître ou la démontrer avant d’appliquer le rapport.
\label{fig-cor-15-10}\figureGeo[corrige]{cos-10}
\end{corrige}\end{blocLivre}
\begin{blocLivre}{cor-15-11}\begin{corrige}{11}{Inclinaison d’une échelle}\demarche{Combiner cosinus pour l’angle et Pythagore pour la hauteur.}
Pour l’angle au sol, adjacent = 1,5 et hypoténuse = 4. L’angle vaut $\arccos(1{,}5/4)\approx\res{68^\circ}$. Pythagore donne $h=\sqrt{4^2-1{,}5^2}\approx\res{3{,}71\ \mathrm m}$.
\label{fig-cor-15-11}\figureGeo[corrige]{cos-11}
\end{corrige}\end{blocLivre}
\begin{blocLivre}{cor-15-12}\begin{corrige}{12}{Une rampe}\demarche{Identifier projection adjacente et longueur inclinée.}
La rampe est l’hypoténuse du triangle rectangle. $\cos\alpha=5{,}7/6=0{,}95$, donc $\alpha=\arccos(0{,}95)\approx\res{18{,}2^\circ}$. La faible différence entre les deux longueurs correspond à une inclinaison modérée.
\label{fig-cor-15-12}\figureGeo[corrige]{cos-12}
\end{corrige}\end{blocLivre}
\begin{blocLivre}{cor-15-13}\begin{corrige}{13}{Un chemin incliné}\demarche{Multiplier l’hypoténuse par le cosinus de l’inclinaison.}
La projection est le côté adjacent, le chemin l’hypoténuse. $d=3{,}2\cos27^\circ\approx\res{2{,}85\ \mathrm{km}}$. Elle est inférieure à la longueur réellement parcourue.
\label{fig-cor-15-13}\figureGeo[corrige]{cos-13}
\end{corrige}\end{blocLivre}
\begin{blocLivre}{cor-15-14}\begin{corrige}{14}{Comparer deux triangles}\demarche{Conserver l’hypoténuse commune dans les deux calculs.}
Les côtés adjacents valent $10\cos30^\circ\approx\res{8{,}66}$ cm et $10\cos60^\circ=\res5$ cm. Pour cette hypoténuse commune, le côté adjacent au plus grand angle est plus court. La comparaison vient des deux calculs, pas de la seule apparence.
\label{fig-cor-15-14}\figureGeo[corrige]{cos-14}
\end{corrige}\end{blocLivre}
\begin{blocLivre}{cor-15-15}\begin{corrige}{15}{Deux outils pour une hauteur}\demarche{Déterminer l’hypoténuse avant d’utiliser Pythagore.}
ABC est rectangle en A. $\cos58^\circ=9/BC$, donc $BC=9/\cos58^\circ\approx\res{17{,}0}$ m. Avec sa valeur non arrondie, $AC=\sqrt{BC^2-9^2}\approx\res{14{,}4}$ m. Le résultat intermédiaire complet évite d’accumuler les erreurs d’arrondi.
\label{fig-cor-15-15}\figureGeo[corrige]{cos-15}
\end{corrige}\end{blocLivre}
